\documentclass[12pt]{article}
\usepackage[utf8]{inputenc}
\usepackage[T1]{fontenc}
\usepackage{hyperref}
\title{Mawa Language Resources for K-12 Curriculum Development}
\author{Compiled by Claude Sonnet 5.5 (claude-sonnet-5-5) for Andr\'as Kornai}
\date{}
\begin{document}
\maketitle

\section{Introduction}
Mawa (ISO 639-3: wma; Glottocode mawa1236) is listed by Glottolog as unattested, in Nigeria: \url{https://glottolog.org/resource/languoid/id/mawa1236.bigmap.html}. Localizely labels the code as extinct: \url{https://localizely.com/language-code/wma}. The searches did not confirm the retirement of the code, and found no evidence that the language exists. It is not to be confused with Mawa (Chad), a different language with code mcw.

\section{Primary Educational Resources}
None exist as far as the searches showed.

\section{Linguistic and Reference Materials}
None found. Only code registry entries exist.

\section{Religious and Community Resources}
None found.

\section{Digital and Audio Resources}
None found.

\section{Curriculum Development Context}
No curriculum can be built for this language. There is no evidence of speakers, a community or any documentation.
\end{document}